\chapter{Bott-сравнение обменного вещественного класса Тёплица}

Ретроспективный аудит свёл вопрос к одному отображению:
\begin{equation}
 c_6:KO_6(\mathbb C_{\mathbb R})
 \longrightarrow
 K_6(\mathbb C\otimes_{\mathbb R}\mathbb C)
 \simeq K_0(\mathbb C\oplus\mathbb C).
 \label{eq:v5-bott-comparison-target}
\end{equation}
Если генератор переходит в антидиагональную пару $(-1,+1)$, то уже
построенные индексы $(-15,+15)$ имеют единственный вещественный прообраз
пятнадцать. Если образ диагонален, ветвь закрывается.

\section{Явная таблица из вещественной K-теории}

В примере 10.2 монографии Бурсема--Шокета
\path{arXiv:2407.05880v4} вычислена объединённая $K$-теория комплексных
чисел, рассматриваемых как вещественная $C^*$-алгебра. Комплексификация
расщепляется:
\begin{equation}
 \mathbb C_{\mathbb R}\otimes_{\mathbb R}\mathbb C
 \simeq\mathbb C\oplus\mathbb C.
 \label{eq:v5-bott-c-complexification}
\end{equation}
В чётных степенях таблица имеет вид
\begin{align}
 c_0(n)&=(n,n),&
 c_2(n)&=(-n,n),\
 c_4(n)&=(n,n),&
 \boxed{c_6(n)=(-n,n).}
 \label{eq:v5-bott-c-even-maps}
\end{align}
Тот же источник даёт в степени шесть операции
\begin{equation}
 r_6(a,b)=-a+b,
 \qquad
 \psi_6(a,b)=(-b,-a),
 \label{eq:v5-bott-r-psi-six}
\end{equation}
где $r$ забывает комплексную структуру, а $\psi$ является комплексным
сопряжением. Общая KK-теоретическая конструкция отображений $c$ и $r$ и
тождество $r\circ c=2$ доказаны Шиком; см.
\path{arXiv:math/0311295}.

Подстановка генератора даёт
\begin{equation}
 c_6(1)=(-1,+1),
 \qquad
 r_6c_6(1)=2,
 \qquad
 \psi_6c_6(1)=c_6(1).
 \label{eq:v5-bott-generator-checks}
\end{equation}
Следовательно, $c_6:\mathbb Z\to\mathbb Z^2$ инъективно и его образ
равен антидиагональной подгруппе
\begin{equation}
 \{(-n,n):n\in\mathbb Z\}.
 \label{eq:v5-bott-antidiagonal-image}
\end{equation}

Знак не является произвольной вставкой. В комплексной $K$-теории
сопряжение меняет знак Bott-элемента степени два; в степени шесть действует
его третья степень, поэтому две компоненты имеют противоположные знаки.

\section{Применение к коэффициентному проектору}

Построенная обменная пара имеет комплексные компоненты
\begin{equation}
 \operatorname{ind}T_+=-15,
 \qquad
 \operatorname{ind}T_-=+15.
 \label{eq:v5-bott-existing-pair}
\end{equation}
По формуле \eqref{eq:v5-bott-c-even-maps}
\begin{equation}
 c_6(15)=(-15,+15).
 \label{eq:v5-bott-fifteen-image}
\end{equation}
Это в точности пара \eqref{eq:v5-bott-existing-pair}, включая порядок
ориентаций $T_+=S\otimes q_0$ и
$T_-=S^*\otimes\overline{q_0}$.

Поскольку $c_6$ инъективно, прообраз единственен.

\begin{theorem}[Вещественный индекс коэффициентного цикла]
Сбалансированный обменный цикл Тёплица представляет класс
\begin{equation}
 [T_{\mathbb R}]=15
 \quad\text{в}\quad
 KO_6\bigl(M_{105}(\mathbb C)_{\mathbb R}\bigr)
 \simeq\mathbb Z,
 \label{eq:v5-bott-real-class-fifteen}
\end{equation}
с точностью до одновременного обращения выбранной хопфовой ориентации.
Его комплексификация есть пара индексов $(-15,+15)$.
\end{theorem}

Нормированное следовое чтение класса равно
\begin{equation}
 \frac{|[T_{\mathbb R}]|}{105}
 =\frac{15}{105}
 =\frac17.
 \label{eq:v5-bott-real-one-seventh}
\end{equation}
Тем самым два ранее разделённых утверждения теперь соединены:
целочисленный вещественный класс равен пятнадцати, а $1/7$ является его
нормированным коэффициентным весом.

\section{Что именно теперь закрыто}

Явная запись матриц $Cl_{0,6}$ больше не является логическим условием
классификации: степень шесть, группа, комплексification и инъективный образ
уже вычислены стандартной Bott-таблицей, а проектный цикл имеет ровно этот
образ. Clifford-линейная модель остаётся полезным конкретным
представителем, но не может изменить полученный класс.

Однако теорема \eqref{eq:v5-bott-real-class-fifteen} относится к
ограниченному аналитическому циклу Тёплица. Она не выводит из конечного
$M_{35}$:
\begin{enumerate}
 \item неограниченный самосопряжённый родительский оператор;
 \item поляризацию Харди как внутреннее следствие геометрии;
 \item пространственную локализацию и конечную энергию;
 \item массу, масштаб или наблюдаемую частицу.
\end{enumerate}

\begin{nogocriterion}[Запрет физического чтения одного K-класса]
Число пятнадцать и вес $1/7$ теперь являются доказанными топологическими
данными обменного цикла. Они не считаются массой, числом поколений,
вероятностью или энергией без отдельного родительского оператора и
вариационного действия.
\end{nogocriterion}

Следующий конструктивный вопрос --- можно ли получить тот же ограниченный
цикл как ограниченное преобразование неограниченного вещественного
оператора, выведенного из хопфово-моритовского соответствия, а не
добавленного внешним пространством Харди.

\begin{protocol}
\begin{enumerate}
 \item группа $KO_6(\mathbb C_{\mathbb R})$: \textbf{$\mathbb Z$};
 \item комплексification коэффициентной алгебры:
 \textbf{$\mathbb C\oplus\mathbb C$};
 \item образ генератора: \textbf{$(-1,+1)$};
 \item инъективность $c_6$: \textbf{выполнена};
 \item проверка $r_6c_6=2$: \textbf{выполнена};
 \item проверка сопряжения $\psi_6$: \textbf{выполнена};
 \item образ класса пятнадцать: \textbf{$(-15,+15)$};
 \item совпадение с парой Тёплица: \textbf{точное};
 \item вещественный класс: \textbf{$15$ с точностью до ориентации};
 \item нормированный вес: \textbf{$1/7$};
 \item явные Clifford-матрицы как условие классификации:
 \textbf{больше не требуются};
 \item неограниченный родительский цикл: \textbf{не построен};
 \item физическое замыкание: \textbf{не объявлено};
 \item следующий гейт:
 \textbf{неограниченный вещественный родительский цикл Тёплица}.
\end{enumerate}
\end{protocol}

Машинный сертификат сохранён в
\path{s2t_v5_real_toeplitz_bott_comparison_map_gate_results.json}.